\subsection*{Maximum Flow and Minimum Cut Problems}

We now define the \emph{maximum-flow problem}: given network $(G= (V, E), s, t, c(\cdot))$,
to find an $s$-$t$ flow $f$ such that $|f|$ is maximized.

See Figure~\ref{fig:maxflow} for an example, which shows a flow with value of 9. 
This flow has a larger value than the flow in the previous figure (on the same network). 
Can you play with this example to try to obtain a flow with even larger value?
Is this flow~(of value 9) the maximum flow of the network?
If it is, how can we verify it is one maximum flow?

%Playing with the instance you may find
%a flow $f$ with $|f| = 9$, shown with red numbers in the figure.
%It seems that this flow $f$ is a maximum flow.
%Here is an argument that this flow $f$ is indeed a maximum flow.
%Consider the \emph{cut} shown with blue curve. Notice that
%$s$ and $t$ are on different sides of this cut. Therefore, all flow
%from $s$ to $t$ must cross this cut. Notice also that there are three
%cut-edges~(i.e., edges cross the border from $s$-side to $t$-side)
%and the \emph{total capacities} of these 3 edges is $3 + 2 + 4 = 9$.
%This gives an \emph{upper bound} of the maximum flow, which is 9.
%And, we have found a flow $f$ with $|f| = 9$. Therefore, $f$
%is a maximum flow.

\begin{figure}[h]
\centering{\input{maxflow}}
\caption{A flow with value 9.} %~(same network with the right panel of Figure 1.}
\label{fig:maxflow}
\end{figure}


%\subsection*{Minimum-Cut}

%A \emph{cut} of a graph $G = (V, E)$ is defined
%as a \emph{partition} of the vertices, formally written as $(S, V\setminus S)$ where $S \subset V$.
We answer above question using $s$-$t$ cut.
Given a network $G=(V, E)$ with source $s\in V$, sink $t\in V$, and capacity $c(\cdot)$,
an \emph{$s$-$t$ cut} of the network is a pair $(S, T)$, where $S\subset V$, $T = V\setminus S$,
and that $s\in S$ and $t\in T$. In short, an $s$-$t$ cut of a network is a partition
of the vertices that separates the source and the sink.

An $s$-$t$ cut $(S, T)$ of a network also partitions all
edges in the graph into four disjoint subsets, described below. 
We define the first category, i.e., $E(S, T) := \{(u,v)\in E\mid u\in S, v\in T\}$,
as the \emph{cut-edges} w.r.t.\ the $s$-$t$ cut $(S, T)$.
\vspace*{-\topsep}
\begin{enumerate}
\item Edges that span the cut and point from $s$-side to $t$-side, formally $E(S, T) := \{(u,v)\in E\mid u\in S, v\in T\}$;
\item Edges that span the cut and point from $t$-side to $s$-side, formally $E(T, S) := \{(u,v)\in E\mid u\in T, v\in S\}$;
\item Edges with both end-vertices in $s$-side, formally $E(S, S) := \{(u,v)\in E\mid u\in S, v\in S\}$;
\item Edges with both end-vertices in $t$-side, formally $E(T, T) := \{(u,v)\in E\mid u\in T, v\in T\}$.
\end{enumerate}

Let $(S, T)$ be an $s$-$t$ cut of a network $G = (V, E)$ with source $s$, sink $t$, and capacity $c(e)$ for any $e\in E$.
We define the \emph{capacity} of $(S, T)$, denoted as $c(S, T)$, as the sum of the capacities of its cut-edges. Formally
as $c(S, T) := \sum_{e\in E(S, T)} c(e)$. See Figure~\ref{fig:flow-cut}.

\begin{figure}[h]
\centering{\input{cut}}
\caption{Four different $s$-$t$ cuts and their capacities of the same network.}
\label{fig:flow-cut}
\end{figure}

We now introduce the so-called \emph{minimum-cut problem}: given network $G = (V, E)$ with source $s$, sink $t$ and capacity $c(e)$ for any $e\in E$,
   we seek an $s$-$t$ cut $(S, T)$ such that its capacity $c(S, T)$ is minimized.

Note that, the definition of $s$-$t$ cut, capacity of an $s$-$t$ cut, and above minimum-cut problem,
are completely independent of flow. Note too that, the input of the maximum-flow problem and the minimum-cut
problem is the same: a network~(consisting of a directed graph, source and sink, and edge capacities).

Although the maximum-flow problem and the minimum-cut problem are defined independently, they
can be solved by the same algorithm~(for example, the Ford-Fulkson algorithm)
and are connected by an elegant theorem~(i.e., the max-flow min-cut theorem).
Below, we first show half of this theorem.
Later on we will introduce the Ford-Fulkson algorithm
and use it to prove the other other half of the theorem.

%\subsection*{Maximum-Flow v.s.\ Minimum-Cut, One-Side Connection}

We now show that, the capacity of \emph{any} $s$-$t$ cut of a network
gives an \emph{upper bound} of the value of \emph{any} $s$-$t$ flow of the same network, 
formally stated below.
\begin{claim} \label{claim1}
Let $(G=(V, E), s, t, c(\cdot))$ be a network.
Let $f$ be an arbitrary $s$-$t$ flow and let $(S, T)$ be an arbitrary $s$-$t$ cut, of this network.
We have $|f| \le c(S, T)$.
\end{claim}

Above claim is intuitive to understand: an $s$-$t$ flow $f$ transfers $|f|$ units of flow from source $s$ to sink $t$;
since an $s$-$t$ cut $(S, T)$ separates $s$ and $t$, so the flow, all generated from $s$, must cross the cut in order to reach $t$,
and such ``crossing'' must use the cut-edges.
Hence, the total amount of the flow that can be transferred, i.e., the value of $f$, is limitted by
the total capacities of the cut-edges, i.e., the capacity of the $s$-$t$ cut.
See Figure~\ref{fig:claim}.

\begin{figure}[h]
\centering{\input{claim}}
\caption{An example showing Claim~\ref{claim1} and Fact~\ref{fact1}.}
\label{fig:claim}
\end{figure}


In order to formally prove above claim, we first show that, the value of an $s$-$t$ flow $f$
can also be calculated by examining the flow carried by the ``spanning edges'' of any $s$-$t$ cut.
Formally,
\begin{fact} \label{fact1}
Let $f$ be an arbitrary $s$-$t$ flow and let $(S, T)$ be an arbitrary $s$-$t$ cut, of network $(G=(V, E), s, t, c(\cdot))$.
We have $|f| = \sum_{e\in E(S, T)} f(e) - \sum_{e\in E(T, S)} f(e)$.
\end{fact}

Above fact is also intuitive to understand: $S$ is a subset that includes $s$ but excludes $t$.
See Figure~\ref{fig:claim}.
When thinking $S$ as a whole, 
the total amount of flow generated by $S$ is equal to $|f|$, as only $s$ generates flow.
On the other hand, this amount of flow emitted from $S$ 
can also be calcuated by summing up the flow leaves $S$, i.e., $\sum_{e\in E(S, T)} f(e)$,
and subtracting the amount of the flow that enters back to $S$, i.e., $\sum_{e\in E(T, S)} f(e)$.

We now formally prove Fact~\ref{fact1}. 
%Similar to the proof of Fact~1 of Lecture~36, 
		This proof also simply applies the conservation property.

\emph{Proof of Fact~\ref{fact1}.} According to conservation property, for any $v\in S\setminus \{s\}$,
we have $\sum_{e\in I(v)} f(e) = \sum_{e\in O(v)} f(e)$. 
We sum up both sides over all $v\in S\setminus\{s\}$: 
we have $\sum_{v\in S\setminus\{s\}} \sum_{e\in I(v)} f(e) = \sum_{v\in S\setminus\{s\}} \sum_{e\in O(v)} f(e)$. 
Note that $I(s) = \emptyset$ as $s$ is a source, so the left side of above equation can be rewritten as $\sum_{v\in S} \sum_{e\in I(v)} f(e)$.
Note that $|f| = \sum_{e\in O(s)} f(e)$ by definition, so the right side can be rewritten as $\sum_{v\in S} \sum_{e\in O(v)} f(e) - |f|$.
Combined, we have $|f| = \sum_{v\in S} \sum_{e\in O(v)} f(e) - \sum_{v\in S} \sum_{e\in I(v)} f(e)$. 
Now consider the 4 types of edges partitioned by $(S, T)$ and consider if they are taken into account by the two items
i.e., $\sum_{v\in S} \sum_{e\in O(v)} f(e)$ and $\sum_{v\in S} \sum_{e\in I(v)} f(e)$. 
\vspace*{-\topsep}
\begin{enumerate}
\item All edges in $E(S, T)$ appear in the first item, but not the second one.
\item All edges in $E(T, S)$ appear in the second item, but not the first one.
\item All edges in $E(S, S)$ appear in both items~(so they cancel out).
\item None of the edges in $E(T, T)$ appear in either item.
\end{enumerate}
Therefore, we have $|f| = \sum_{e\in E(S, T)} f(e) - \sum_{e\in E(T, S)} f(e)$. \qed

We now formally prove Claim~\ref{claim1}. The proof simply combines the capacity condition (and the non-negative property) of an $s$-$t$ flow,
Fact~\ref{fact1}, and the definition of the capacity of an $s$-$t$ cut.

\emph{Proof of Claim~\ref{claim1}.} Since $0\le f(e) \le c(e)$ for any $e\in E$,
based on Fact~\ref{fact1}, we can write
$|f| = \sum_{e\in E(S, T)} f(e) - \sum_{e\in E(T, S)} f(e)
\le \sum_{e\in E(S, T)} c(e) - 0 = c(S, T)$. \qed


Claim~\ref{claim1} states that capacity of \emph{any} $s$-$t$ cut
is larger or equal to the value of \emph{any} $s$-$t$ flow.
Consequently, if there exists an $s$-$t$ flow $f$ and an $s$-$t$ cut $(S,T)$
satisfies that $|f| = c(S, T)$, then this $f$ must be a maximum-flow,
and this $(S, T)$ must be a minimum-cut.
An illustration of this property is given in Figure~\ref{fig:flow-meet},
where each red point represents the value of an $s$-$t$ flow
while each blue point represents the capacity of an $s$-$t$ cut.
So all red points lie to the left side of all blue points,
and if a red point and a blue point meet, they are maximum-flow and minimum-cut.

\begin{figure}[t]
\centering{\input{flow-meet}}
\caption{Illustration of the relationship between capacity of any $s$-$t$ cut
and value of \emph{any} $s$-$t$ flow.}
\label{fig:flow-meet}
\end{figure}


\begin{figure}[h]
\centering{\input{maxflow-mincut}}
\caption{A maximum-flow $f$ and a minimum cut $(S, T)$, verified by $|f| = c(S,T)$.}
\label{fig:maxflow-mincut}
\end{figure}

In above example, we see a flow $f$ with $|f| = 9$ and an $s$-$t$ cut $(S, T)$ with $c(S, T) = 9$.
So both are optimal. See Figure~\ref{fig:maxflow-mincut}.

%At the end of previous Lecture we gave an example~(network) for which
%there exists a flow $f$ and an $s$-$t$ cut such that $|f| = c(S, T)$.
%Therefore, both $f$ and $(S,T)$ are optimal.
Is this the case for \emph{any} network?
That is, for an arbitrary network,
does there \emph{always} exist $f$ and $(S,T)$ such that $|f| = c(S,T)$?
The answer is yes! This elegant, surprising result is referred to the max-flow min-cut theorem.
The proof is constructive: we will design an algorithm~(i.e., the Ford-Fulkson algorithm)
that actually always finds an $s$-$t$ flow $f$ and an $s$-$t$ cut $(S, T)$ that satisfies $|f| = c(S, T)$.

